\documentclass[10pt]{article}
\usepackage[a4paper,margin=27mm]{geometry}
\usepackage[T1]{fontenc}
\usepackage{lmodern}
\usepackage{amsmath,amssymb,amsthm,mathtools}
\usepackage{microtype}
\usepackage[hidelinks]{hyperref}
\hypersetup{pdftitle={A two-tail obstruction to ambient conjugation in Dlab groups},pdfauthor={Alec Kriebel}}
\usepackage{enumitem}
\setlist{nosep}
\newtheorem{theorem}{Theorem}
\newtheorem{lemma}[theorem]{Lemma}
\newtheorem{proposition}[theorem]{Proposition}
\newtheorem{corollary}[theorem]{Corollary}
\theoremstyle{definition}
\newtheorem{definition}[theorem]{Definition}
\newtheorem{remark}[theorem]{Remark}
\newcommand{\R}{\mathbb R}
\newcommand{\Z}{\mathbb Z}
\newcommand{\id}{\mathrm{id}}
\newcommand{\supp}{\operatorname{supp}}
\newcommand{\Fix}{\operatorname{Fix}}
\newcommand{\Homeo}{\operatorname{Homeo}}
\newcommand{\Aff}{\operatorname{Aff}}
\newcommand{\AI}{\mathcal A_I}
\newcommand{\AR}{\mathcal A_{\R}}
\newcommand{\GL}{G(L)}
\title{A two-tail obstruction to ambient conjugation in Dlab groups}
\author{Alec Kriebel\\
\small ORCID: \href{https://orcid.org/0009-0001-9320-500X}{0009-0001-9320-500X}}
\date{17 October 2026\\[3pt]
\small Version 1.0.0}
\begin{document}
\maketitle

\begin{abstract}
Let $G=D_{\langle2\rangle}([0,1])$ be the compactly supported,
locally right $2^{\Z}$-affine Dlab group. We construct an automorphism
preserving both Dlab orders and give a proof that it cannot be implemented
after any embedding in either of two explicitly defined groups: all locally
right affine increasing interval homeomorphisms, and all locally right
affine increasing real homeomorphisms that are affine on an initial ray.
The ambient slopes are unrestricted, and the embedding need not preserve
an order. Perfectness separates the supports of commuting local groups;
a two-tail construction handles both possible orientations of their least
action components. This extends a one-tail mechanism used in prior work.
The definitions and affine-germ decomposition of the standard Dlab
families place them in these ambient groups. Consequently the construction
gives an affirmative answer to the corrected Kourovka Problem 21.149,
including arbitrary positive-real slope groups and the full and reflected
extended-real families.
\end{abstract}

\section{The precise claim and its context}

For an interval $X$, let $\Homeo^+(X)$ denote its increasing
homeomorphisms, with ordinary composition of functions. For $f$ in this
group, $\supp(f)=\{x:f(x)\ne x\}$ is its open support. An open support
component is called a support component. For intervals $U,V$, write $U<V$
if every point of $U$ precedes every point of $V$.

\begin{definition}
A map is \emph{locally right affine} if, at every finite point $a$ other
than a right endpoint, some $\epsilon>0$ and $s>0$ satisfy
\[
 f(x)=f(a)+s(x-a)\qquad(a\le x\le a+\epsilon).
\]
Let $\AI\le\Homeo^+([0,1])$ consist of all such maps, including this
condition at $a=0$. Let $\AR\le\Homeo^+(\R)$ consist of all such maps
which, in addition, agree with $x\mapsto sx+t$, for some $s>0$, on an
initial ray. No restriction is placed on the positive slopes.
\end{definition}

These are groups: right affine germs are preserved by composition and
inversion, as are eventual affine germs at $-\infty$. Neither definition
requires finitely many affine pieces on a compact interval; accumulation
of pieces from the left is permitted. Put
\[
 G=\{f\in\AI:\text{all right slopes lie in }2^{\Z},\quad
 f=\id\text{ near }0\text{ and }1\}.
\]
This is Dlab's compact-support group. For $L=(a,b)$ with
$0<a<b<1$, put
\[
 G(L)=\{f\in G:\overline{\supp(f)}\text{ is compactly contained in }L\}.
\]

\begin{theorem}[Explicit ambient theorem]\label{thm:main}
There is an automorphism $\alpha$ of $G$ preserving both Dlab bi-orders
such that, for $A=\AI$ or $A=\AR$, every injective homomorphism
$e:G\longrightarrow A$ and every $u\in A$ satisfy
\[
 e(\alpha(f))\ne u^{-1}e(f)u\quad\text{for some }f\in G.
\]
Consequently the same assertion holds for every subgroup of either
explicit ambient group. No order on $A$ or order preservation by $e$ is
assumed.
\end{theorem}

Problem 21.149 in the 21st \emph{Kourovka Notebook}, updated
6 October 2026, printed p.~199, is attributed to V.~M.~Kopytov and
N.~Ya.~Medvedev and communicated by A.~V.~Zenkov~\cite{KN}.
It asks for an order automorphism of a Dlab group whose realization by
conjugation remains impossible in a possibly larger Dlab group. The
earlier question asked only for an automorphism which is not inner;
van Doorn, Judin, Monticone and Morrison discuss that distinction and
the clarification in Appendix~A of~\cite{VDJMM}.

The associated Lean formalization explicitly labels its final theorem as
solving the original formulation.\footnote{\href{https://github.com/pitmonticone/Kourovka/blob/64a4fe59b1df1893a0a1a68641d6806ccaaa7150/Kourovka/Problem_21_149.lean\#L768-L775}{Pinned Lean formalization},
lines 768--775; accessed 17 October 2026.}
Its conjugator is the finite piecewise affine map
\[
 h_0(x)=2x\quad(0\le x\le1/3),\qquad
 h_0(x)=(x+1)/2\quad(1/3\le x\le1).
\]
This map belongs to the full interval group $\overline D_{\langle2\rangle}(I)$,
although it is outside the compact-support source $G$. Thus that
particular non-inner automorphism is implemented in a larger Dlab group,
as allowed by the original assertion but excluded by the corrected target.

Gong, Yang and Zeng~\cite[Theorem~1.1]{GYZ} already give an ambient
nonconjugation result for six rank-one Dlab families and order-preserving
embeddings. Their proof uses a normalizer with support components
descending to one endpoint. They explicitly exclude higher-rank ambient
slopes and the reflected extended-real family from their theorem. We
credit that one-tail support obstruction. The additional argument here
uses perfect centralizers instead of ambient order classification, and
uses both tails to accommodate an embedding that reverses the order
of least local action components. The conclusion concerns the stronger ambient quantifiers, rather than
a new claim about the already known one-tail mechanism.

\section{Germs and perfect centralizers}

The external group theorem needed for the support obstruction is Dlab's
simplicity theorem~\cite[Theorem~4.3]{Dlab}. It states that $G$ is
algebraically simple. The group is nontrivial, for it contains finite
two-slope bumps. An increasing homeomorphism of an interval has no
nontrivial finite order: a moved point has strictly monotone iterates.
Thus $G$ cannot be an abelian simple group, and is nonabelian and
perfect. An affine change from $[0,1]$ to $[a,b]$, followed by extension
by the identity, identifies $G$ with $G(L)$. In particular every $G(L)$
is also nonabelian simple and perfect.

\begin{lemma}\label{lem:wellordered}
For every $q\in\AI$ or $q\in\AR$, its support components are
well ordered from left to right. If $q$ belongs to a perfect subgroup
of the ambient group, every support component has a finite left endpoint.
\end{lemma}

\begin{proof}
First consider a nonempty family of components with finite left
endpoints whose infimum $a$ is finite. If no component has left endpoint
$a$, choose distinct such endpoints $a_n\downarrow a$. They are fixed
points of $q$, so continuity gives $q(a)=a$. A right affine germ at $a$
then has slope $1$, because it contains the fixed points $a_n$ for
arbitrarily large $n$. It is therefore the identity on a right
neighborhood of $a$. No left endpoint of a nonempty open support
component can lie in that neighborhood, a contradiction. Thus the
family has a first component. On $[0,1]$ its infimum is always finite
and less than $1$, so this proves the assertion there.

On $\R$, if $q$ is the identity on an initial ray, all support components
have left endpoints bounded below and the argument applies. Otherwise
its eventual nonidentity affine formula moves every point sufficiently
far to the left. There is a single initial support component, with left
endpoint $-\infty$, and all other component endpoints are bounded
below. A family containing that component has it as first member; a
family excluding it is handled by the finite-infimum argument.

For $\AI$, the right slope at $0$ is a homomorphism to the abelian
group $\R_{>0}^{\times}$. It is trivial on any perfect subgroup, and a
map fixing $0$ with slope $1$ is the identity on an initial interval.
For $\AR$, the eventual affine formula defines a homomorphism into
$\Aff_+(\R)$. This group is metabelian: its slope quotient and
translation kernel are abelian. A perfect subgroup has trivial image
in a solvable group, so every element of the perfect subgroup is the
identity on an initial ray. These neighborhoods can depend on the
element; no common neighborhood is asserted.
\end{proof}

\begin{lemma}[Perfect centralizers]\label{lem:centralizer}
Let $q$ be an ambient element all of whose support components have
finite left endpoints. If a perfect subgroup $P$ of the ambient group
centralizes $q$, then every element of $P$ fixes $\supp(q)$ pointwise.
\end{lemma}

\begin{proof}
Each $p\in P$ permutes the support components of $q$ increasingly.
An increasing bijection of a well-ordered set is the identity, by
transfinite induction: if every predecessor of an element is fixed,
its image cannot be an earlier element, already having its own
preimage, and cannot be a later element, since its own preimage would
then lie later and violate monotonicity.
Lemma~\ref{lem:wellordered} therefore implies that $P$ preserves
each component $J=(a,b)$ setwise. Here $a$ is finite, so $P$ fixes $a$.
The right slope at $a$ is a homomorphism $P\to\R_{>0}^{\times}$;
perfectness makes it trivial. Every $p\in P$ is consequently the
identity on a right neighborhood of $a$.

The map $q$ has no fixed point in $J$, and either $q(x)>x$ throughout
$J$ or $q(x)<x$ throughout $J$. Accordingly either $q^{-n}(x)$ or
$q^n(x)$ tends to $a$ for each $x\in J$: a larger finite limit would
be a fixed point by continuity. For $n$ sufficiently large this
iterate lies in the neighborhood fixed by $p$. Commutation with $q$
then implies $p(x)=x$. This proves the assertion for every $J$ and $p$.
\end{proof}

For a subgroup $B$ define its action support by
\[
 S(B)=X\setminus\Fix(B)=\bigcup_{b\in B}\supp(b).
\]
An \emph{action component} is a component of this open set. Finite
endpoints of an action component are common fixed points of $B$;
hence every element of $B$ preserves every action component setwise.

\begin{corollary}\label{cor:disjoint}
For every embedding $e:G\to A$ as in Theorem~\ref{thm:main}, disjoint
source intervals $L,M$ satisfy
$S(e(G(L)))\cap S(e(G(M)))=\varnothing$.
\end{corollary}

\begin{proof}
The groups $G(L)$ and $G(M)$ commute since their supports are disjoint.
For each $q\in e(G(L))$, its support components have finite left
endpoints by perfectness and Lemma~\ref{lem:wellordered}.
The perfect group $e(G(M))$ centralizes $q$; apply
Lemma~\ref{lem:centralizer} and take the union over $q$.
\end{proof}

\section{Least components and the orientation dichotomy}

Fix an embedding $e:G\to A$. Its action on every action component
$J$ is faithful. Indeed, the restriction kernel is normal in $G$,
and it is proper because $J$ is an action component; simplicity
makes the kernel trivial. Fix $f\ne1$ in $G$. Every such $J$ contains
a support component of $e(f)$, since faithfulness makes its restriction
nonidentity and a support component cannot cross a common fixed
endpoint of $J$. The least support component of $e(f)$ therefore lies
in the least action component of $e(G)$; denote the latter by $J_0$.
To justify minimality explicitly, an earlier action component would
contain an even earlier support component of $e(f)$.

For each source interval $L$, the action of $e(G(L))$ on $J_0$ is
nontrivial, since the action of $e(G)$ there is faithful. Every local
action component within $J_0$ carries a faithful action of $G(L)$,
by its simplicity. Applying the preceding fixed-nonidentity-element
argument to $G(L)$ gives a least such component, denoted $K(L)$.
These are also global action components of $e(G(L))$: the boundary
of $J_0$ is fixed by $e(G)$ and thus cannot be crossed by a local
action component. Corollary~\ref{cor:disjoint} shows that $K(L)$ and
$K(M)$ are disjoint whenever $L$ and $M$ are disjoint.

Conjugation by any $g\in G$ maps $G(L)$ onto $G(gL)$. Moreover
$e(g)$ preserves $J_0$ and sends least components to least components.
Consequently
\begin{equation}\label{eq:equivariant}
 K(gL)=e(g)K(L)\qquad(g\in G).
\end{equation}

\begin{lemma}\label{lem:orientation}
Exactly one of the following holds for every pair of relatively
compact source intervals $L<M$ with a positive gap:
\[
 K(L)<K(M);\qquad K(M)<K(L).
\]
\end{lemma}

\begin{proof}
We first verify finite-tuple transitivity. Given two increasing
tuples of distinct interior points of equal length, choose common
outer endpoints $c,d\in(0,1)$ containing both tuples. Prescribe the
tuple map and fix $c,d$. On any successive source segment of length
$\ell$ and target segment of length $r\ell$, choose $s=2^k$ with
$s\le r\le2s$. Use slope $s$ for length $\ell(2s-r)/s$ and slope
$2s$ for length $\ell(r-s)/s$. These lengths are nonnegative, sum
to $\ell$, and give target length $r\ell$. Omit zero-length pieces.
The resulting finite piecewise affine map, extended by the identity
outside $[c,d]$, belongs to $G$ and realizes the prescribed tuple map.

Apply this construction to the four ordered endpoints of any two
positively separated interval pairs. It maps one pair to the other.
Equation~\eqref{eq:equivariant} and the fact that $e(g)$ increases
show that the relative order of the two $K$-intervals is the same
for every pair. The intervals are disjoint and nonempty, so exactly
one of the stated alternatives holds.
\end{proof}

\section{The two-tail automorphism and the Dlab orders}

For $m\ge0$, set
\[
 b_m=2^{-2m-2},\quad a_m=b_m/2,\quad
 C_m=(a_m,b_m),\quad D_m=(1-b_m,1-a_m).
\]
Then $C_{m+1}<C_m$ and $D_m<D_{m+1}$, each with a positive gap.
On the closure $[a,b]$ of each $C_m$ and $D_m$ define
\begin{equation}\label{eq:bump}
 h(x)=
 \begin{cases}
 a+2(x-a),&a\le x\le a+(b-a)/3,\\
 b+(x-b)/2,&a+(b-a)/3\le x\le b,
 \end{cases}
\end{equation}
and put $h(x)=x$ elsewhere. The two formulas agree at the breakpoint,
fix $a,b$, and satisfy $h(x)>x$ for $a<x<b$. They define increasing
homeomorphisms of each closed bump interval. Because their lengths
tend to zero at both endpoints, their union with the identity is
continuous also at $0,1$, and hence is an increasing homeomorphism
of $[0,1]$. Every compact subinterval of $(0,1)$ meets only finitely
many bumps, so both $h$ and $h^{-1}$ are finite piecewise affine
there, with slopes in $2^{\Z}$.

Define
\begin{equation}\label{eq:alpha}
 \alpha(f)=h^{-1}fh.
\end{equation}
For $f\in G$, its support closure is compactly contained in $(0,1)$,
and the support of $\alpha(f)$ is $h^{-1}\supp(f)$. On the interior,
composition and inversion preserve locally right $2^{\Z}$-affinity.
Near the endpoints the conjugate is the identity. Thus $\alpha(f)$
belongs to $G$; the same argument for $h^{-1}$ proves that $\alpha$
is an automorphism.

For completeness, the orders preserved by $\alpha$ can be specified
without choosing an ambient order. For $1\ne f\in G$, let $a(f)$
be the left endpoint of its first support component, and let $s(f)$
be its right slope there. Existence follows from
Lemma~\ref{lem:wellordered}. We have $0<a(f)<1$ and $s(f)\ne1$;
a slope of $1$ at a fixed point would give an identity right germ.
The condition $s(f)>1$ defines one positive cone. To check closure
under products of positive elements, the earlier first support point
survives unchanged if the two points differ. If they coincide, the
right slopes multiply and their product exceeds $1$. Inversion
reciprocates the slope, and conjugation transports the first support
point increasingly and leaves its slope unchanged, by cancellation
of the conjugator's right slope. This gives a bi-order; the inverse
cone gives its reverse. Dlab's order theorem
\cite[Theorem~4.1]{Dlab} identifies these as the two Dlab orders for
$2^{\Z}$.

The first support point of $h^{-1}fh$ is $h^{-1}(a(f))$. If the
right slope of $h$ there is $r$, the conjugate's first right slope is
$r^{-1}s(f)r=s(f)$. All these germs are at interior points, where
$h$ has the required affine germs. Thus $\alpha$ preserves both
orders, even though $h$ itself has no affine right germ at $0$.

For each $L=C_m$ or $L=D_m$, we have $h(L)=L$, so
\begin{equation}\label{eq:localalpha}
 \alpha(G(L))=G(L),\qquad \alpha|_{G(L)}\ne\id.
\end{equation}
For the second assertion, choose $x\in L$ moved by $h$ and a small
open interval $V$ compactly contained in $L$ with
$V\cap h^{-1}(V)=\varnothing$. A nonidentity finite bump supported
compactly in $V$, obtained by rescaling~\eqref{eq:bump}, belongs to
$G(L)$, whereas its conjugate is supported in $h^{-1}(V)$.
The two elements are distinct.

\section{Proof of the ambient theorem}

\begin{proof}[Proof of Theorem~\ref{thm:main}]
The automorphism and its order preservation were established above.
Suppose, contrary to the remaining assertion, that
\begin{equation}\label{eq:implement}
 e(\alpha(f))=u^{-1}e(f)u\qquad\text{for all }f\in G.
\end{equation}
Then $u$ normalizes $e(G)$, permutes its action components
increasingly, and fixes their least member $J_0$ setwise.
For $L=C_m$ or $D_m$, equations~\eqref{eq:localalpha} and
\eqref{eq:implement} imply that $u$ normalizes $e(G(L))$.
It therefore preserves the least local component $K(L)$ within
$J_0$ setwise.

The restriction of $u$ to $K(L)$ is nonidentity. Otherwise
\eqref{eq:implement} would give equal restrictions of
$e(\alpha(f))$ and $e(f)$ to $K(L)$ for every $f\in G(L)$.
The action of $G(L)$ there is faithful, so this would imply
$\alpha(f)=f$ for all $f\in G(L)$, contradicting
\eqref{eq:localalpha}. Choose a point of $K(L)$ moved by $u$.
Its support component $U(L)$ for $u$ lies entirely inside $K(L)$:
finite boundary points of the invariant interval $K(L)$ are fixed
by $u$, and an infinite boundary cannot be crossed.

If the first orientation of Lemma~\ref{lem:orientation} holds, then
\[
 U(C_{m+1})<U(C_m)\qquad(m\ge0).
\]
If the second orientation holds, then instead
\[
 U(D_{m+1})<U(D_m)\qquad(m\ge0).
\]
In either case these are distinct support components of $u$ forming
a strictly descending sequence. This contradicts
Lemma~\ref{lem:wellordered}, which applies to every ambient element
$u$, whether or not it belongs to a perfect subgroup.
\end{proof}

\section{Why the left-end condition cannot be dropped}

The eventual affine condition at $-\infty$ is a mathematical
hypothesis, not a harmless convention. Here is an exact countermodel
for the larger group $\mathcal B$ of all increasing real homeomorphisms
that are locally right $2^{\Z}$-affine at every finite point and have
support bounded above. No condition at $-\infty$ is imposed.

\begin{proposition}\label{prop:countermodel}
For every automorphism $\beta$ of $G$ there are an embedding
$E:G\to\mathcal B$ and $v\in\mathcal B$ such that
$E(\beta(f))=v^{-1}E(f)v$ for all $f\in G$.
\end{proposition}

\begin{proof}
For $n\in\Z$ put
$I_n=(-2^{-n},-2^{-n-1})$ and let
$\phi_n:[0,1]\to\overline{I_n}$ be the increasing affine bijection.
Define $E(f)$ on $I_n$ by
\[
 E(f)|_{I_n}=\phi_n\,\beta^n(f)\,\phi_n^{-1},
\]
fix every cell endpoint, and make $E(f)$ the identity on $[0,\infty)$.
Each cell is preserved, so continuity holds at $0$ as well as at
the cell endpoints. At a cell endpoint there is an identity right
neighborhood, since $\beta^n(f)$ is the identity near $0$ and $1$.
At $0$ there is an identity right ray. Interior right slopes remain
in $2^{\Z}$ under affine conjugation. Thus $E(f)\in\mathcal B$.
Cellwise composition proves that $E$ is a homomorphism, and its
restriction to $I_0$ proves injectivity.

Put $v(x)=x/2$ for $x<0$ and $v(x)=x$ for $x\ge0$.
This belongs to $\mathcal B$ and satisfies
$v\phi_n=\phi_{n+1}$. The restriction of $v^{-1}E(f)v$ to
$I_n$ is therefore
$\phi_n\beta^{n+1}(f)\phi_n^{-1}$, which equals the restriction
of $E(\beta(f))$. Both maps fix all remaining points.
\end{proof}

In particular it would be incorrect to apply
Theorem~\ref{thm:main} to $\mathcal B$. Its elements can have support
components descending indefinitely toward $-\infty$. We do not
identify $\mathcal B$ with any of the standard real-line Dlab families.

\section{Scope, dependencies, and disclosure}

The explicit theorem includes every interval Dlab subgroup from
Dlab's original definition, with any positive-real slope subgroup
and any group order. The source simplicity and order classification
used here are respectively Theorems~4.3 and~4.1 of~\cite{Dlab};
none of the arguments imposes rank one on an ambient slope group.
There are no computational claims or numerical approximations in
the proof.

For the extended-real families, the definitions and decomposition
in Zenkov--Medvedev~\cite[\S2, Russian pp.~538, 542--543]{ZM}
supply the additional left-end condition. Their full group
$\overline D_H$ consists of locally right $H$-affine maps; the subgroup
$D_{H*}$ fixes an initial ray. In the decomposition proof on p.~542,
every element of $\overline D_H$ is a product of a globally affine
map $x\mapsto rx+t$ and an element of $D_{H*}$. In particular, if
an element is not already in $D_{H*}$, the proof identifies its first
support point as $-\infty$ and gives an affine formula on a right
neighborhood of that point. Such a neighborhood is an initial real
ray. Each full-group element is therefore eventually affine at
$-\infty$, and $\overline D_H\le\AR$. The compact-support,
bounded-below, and reflected bounded-above families are subgroups
of this full group, so they inherit the same conclusion.
This deduction concerns defining affine behavior; it does not apply
a rank-one classification of ambient group orders.
Proposition~\ref{prop:countermodel} explains why retaining that
left-end condition is essential.

\begin{corollary}[Kourovka Notebook 21.149]\label{cor:kourovka}
Equip $G=D_{\langle2\rangle}([0,1])$ with either of its Dlab orders.
The automorphism $\alpha$ of Theorem~\ref{thm:main} cannot be induced
by conjugation after any embedding into any of the standard real Dlab groups
\[
D_H(I),\ D_{H*}(I),\ D_{*H}(I),\ \overline D_H(I),\qquad
D_H,\ D_{H*},\ D_{*H},\ \overline D_H,
\]
for any positive-real slope subgroup $H$. In particular it gives an
affirmative answer to the corrected ambient-conjugation question.
\end{corollary}
\begin{proof}
The four interval groups lie in $\AI$ by their locally right affine
definition. The preceding left-germ argument places all four extended-real
groups in $\AR$. Apply Theorem~\ref{thm:main} to the composite of an
embedding with the appropriate inclusion. The conclusion applies to all
embeddings and hence to all order-preserving embeddings.
\end{proof}

The contribution relative to~\cite{GYZ} is the unrestricted ambient
slope and embedding quantifiers, including the full and reflected
extended-real families. The affine-germ centralizer mechanism is classical
(cf.~\cite[Lemma~3.6 and Theorem~4.6]{Dlab}); the proof above includes
its required one-sided version rather than treating it as a new result.
A literature audit through 17 October 2026 found no prior result covering
these full quantifiers. This is a bounded priority assessment, and no
independent human peer-review or proof-assistant verification is claimed.

Related questions include classifying embeddings beyond their
orientation dichotomy, and characterizing ambient left-germ groups
for which the perfect-centralizer argument survives.

\paragraph{AI use and review status.}
AI tools were used extensively in discovering, developing, writing,
and checking the mathematics, including separate AI research and
adversarial-review agents. These checks do not constitute independent
human peer review. This preprint has not undergone independent human
peer review. No formal proof-assistant verification is claimed.

\begin{thebibliography}{9}

\bibitem{Dlab}
V.~Dlab, \emph{On a family of simple ordered groups},
J.~Austral. Math. Soc. \textbf{8} (1968), no.~3, 591--608.
\href{https://doi.org/10.1017/S1446788700006261}{doi:10.1017/S1446788700006261}.

\bibitem{KN}
E.~I.~Khukhro and V.~D.~Mazurov (eds.),
\emph{Unsolved Problems in Group Theory: The Kourovka Notebook},
21st ed., Sobolev Institute of Mathematics, Novosibirsk, 2026,
update of 6 October 2026, Problem~21.149, p.~199.
\href{https://kourovkanotebookorg.wordpress.com/wp-content/uploads/2026/10/21tkt-1.pdf}{Official edition};
\href{https://kourovkanotebookorg.wordpress.com/wp-content/uploads/2026/10/21upd-1.pdf}{official updates}.

\bibitem{GYZ}
T.~Gong, Y.~Yang and M.~R.~Zeng,
\emph{An order automorphism of a Dlab group not induced by conjugation},
preprint, 16 September 2026, arXiv:2609.18630v1.
\url{https://arxiv.org/abs/2609.18630v1}.

\bibitem{VDJMM}
W.~van Doorn, E.~Judin, P.~Monticone and D.~Morrison,
\emph{On Some Problems from the Kourovka Notebook},
preprint, version~2, 26 July 2026, arXiv:2607.17477v2.
\url{https://arxiv.org/abs/2607.17477v2}.

\bibitem{ZM}
A.~V.~Zenkov and N.~Ya.~Medvedev,
\emph{On Dlab groups}, Algebra and Logic \textbf{38} (1999),
289--298; Russian original: Algebra i Logika \textbf{38} (1999),
no.~5, 531--548.
\href{https://doi.org/10.1007/BF02671746}{doi:10.1007/BF02671746}.

\end{thebibliography}
\end{document}
